% from ega0-19-fr.tex
\oldpage[0\textsubscript{IV-1}]{93}
Prouvons ensuite que $b)$ entraîne $a)$. Le même raisonnement qu'au début de
(19.5.4.2) prouve l'existence d'un $A$-homomorphisme surjectif d'algèbres
$v:F_n\to E_n$ tel que $\gr^j(v)=\varphi_j$ pour tout $j$~; comme $\varphi_j$
est bijectif pour tout $j$ et que les filtrations de $F_n$ et $E_n$ sont
finies, on en conclut que $v$ est \emph{bijectif} (Bourbaki,
\emph{Alg. comm.}, chap.~III, §~2, n\textsuperscript{o}~8, cor.~3 du th.~1).
Soient maintenant $G$ une $A$-algèbre topologique discrète, $\mathfrak R$ un
idéal de carré nul dans $G$, $f:B\to G/\mathfrak R$ un $A$-homomorphisme
continu d'algèbres. Comme $G$ est discrète, il existe un entier $m$ tel que
$f$ s'annule dans $\mathfrak J^m$, donc $f$ se factorise en
$B\to E_n\xrightarrow{f_n}G/\mathfrak R$, où l'on prend $n=2m$. On a donc par
composition un $A$-homomorphisme continu d'algèbres
$r:C\to F_n\xrightarrow{v}E_n\xrightarrow{f_n}G/\mathfrak R$, et comme $C$
est une $A$-algèbre formellement lisse discrète, $r$ se factorise en
$C\xrightarrow{r'}G\to G/\mathfrak R$, de sorte que $G$ est muni par $r'$ d'une
structure de $C$-algèbre. D'autre part, la restriction à
$\mathfrak J/\mathfrak J^2$ de l'homomorphisme
$f'_n:F_n\xrightarrow{v}E_n\xrightarrow{f_n}G/\mathfrak R$ est une application
$C$-linéaire $g:\mathfrak J/\mathfrak J^2\to G/\mathfrak R$. Puisque
$\mathfrak J/\mathfrak J^2$ est un $C$-module projectif, $g$ se factorise en
$\mathfrak J/\mathfrak J^2\xrightarrow{h}G\to G/\mathfrak R$, où $h$ est une
application $C$-linéaire~; par prolongement, on en déduit un homomorphisme de
$C$-algèbres $w:\mathbf S_C^*(\mathfrak J/\mathfrak J^2)\to G$, et par
construction tout élément de degré $m$ a pour image par $w$ un élément de
$\mathfrak R$, donc tout élément de degré $n=2m$ a une image nulle puisque
$\mathfrak R^2=0$~; autrement dit, $w$ se factorise en
$\mathbf S_C^*(\mathfrak J/\mathfrak J^2)\to F_n\xrightarrow{w'}G$. Par
construction, le composé $F_n\xrightarrow{w'}G\xrightarrow{\theta}G/\mathfrak R$
coïncide avec $f'_n$ dans $C$ et dans $\mathfrak J/\mathfrak J^2$, donc est
égal à $f'_n$. Finalement, le composé
$B\to E_n\xrightarrow{v^{-1}}F_n\xrightarrow{w'}G\to G/\mathfrak R$ étant
égal à $f$, on voit que $B$ est une $A$-algèbre formellement lisse.
\end{env}

\begin{env}[19.5.4.4]
\label{0.19.5.4.4-fr}
Il reste à prouver l'équivalence de $a)$ et $c)$. En premier lieu, $c)$ entraîne
$a)$~: en effet, $B'$ est une $C$-algèbre formellement lisse pour la topologie
discrète (19.3.2), donc aussi pour la topologie ${B'}^+$-préadique (19.3.8)~;
comme $C$ est une $A$-algèbre formellement lisse, $B'$ est aussi une
$A$-algèbre formellement lisse (19.3.5, (ii)) et finalement $\widehat{B'}$ est
une $\widehat A$-algèbre formellement lisse (19.3.6), donc $B$ est une
$A$-algèbre formellement lisse (19.3.6). Reste à voir que $a)$ entraîne $c)$.
Notons d'abord que puisque $C$ est une $A$-algèbre formellement lisse,
l'homomorphisme identique $C\to B/\mathfrak J$ se factorise en
$C\to\widehat B\to B/\mathfrak J$, où $C\to\widehat B$ est un
$A$-homomorphisme (19.3.10)~; $\widehat B$, et par suite tous les
$B/\mathfrak J^{n+1}$, sont donc munis de structures de $C$-algèbres. D'autre
part, comme $\mathfrak J/\mathfrak J^2$ est un $C$-module projectif par $b)$,
l'injection canonique $\mathfrak J/\mathfrak J^2\to B/\mathfrak J^2$ permet de
former un système projectif de $C$-homomorphismes
$u_n:\mathfrak J/\mathfrak J^2\to B/\mathfrak J^{n+1}$ pour $n\geq1$, donc
aussi (par la propriété universelle de l'algèbre symétrique) un système
projectif d'homomorphismes de $C$-algèbres
$v_n:\mathbf S_C^*(\mathfrak J/\mathfrak J^2)=B'\to B/\mathfrak J^{n+1}$~;
d'ailleurs il est clair que $v_n$ s'annule dans $({B'}^+)^{n+1}$ et comme
$\gr^0(v_n)=\varphi_0$ et $\gr^1(v_n)=\varphi_1$, on a
$\gr^j(v_n)=\varphi_j$ pour $0\leq j\leq n$. Comme les $\varphi_j$ sont des
isomorphismes par $b)$, et que les filtrations de
$B'/({B'}^+)^{n+1}$ et de $B/\mathfrak J^{n+1}$ sont finies, on en conclut que
$v_n$ est \emph{bijectif} pour

